\chapter{Experimental Design}

\section{Simulation Settings}

The experimental design is intended to make the mechanism comparison
reproducible. Hardware is reported for runtime context, but the
economically relevant controls are the bank population, daily horizon,
common initial-state rules, regulatory threshold, loan cap, random
seeds, and Monte Carlo repetitions. Centralized matching and RFQ are
evaluated under the same settings so that different outcomes do not come
from different balance-sheet calibrations.

\subsection{Computational Environment}

All simulation experiments in the present study are conducted on a
personal computer. The main hardware and system specifications are as
follows:


\begin{table}[htbp]
\centering
\caption{Computational environment}
\label{tab:computational-environment}
\small
\renewcommand{\arraystretch}{1.15}
\begin{tabularx}{0.92\textwidth}{@{}>{\raggedright\arraybackslash}p{0.28\textwidth}X@{}}
\toprule
\textbf{Item} & \textbf{Specification}\\
\midrule
Operating system & Windows 11 Pro\\
Processor & 12th Gen Intel Core i5-12600KF, 3.70 GHz\\
Memory & 32.0 GB RAM\\
Graphics card & NVIDIA GeForce RTX 4060 Ti, 8 GB\\
System type & 64-bit operating system, x64-based processor\\
\bottomrule
\end{tabularx}
\end{table}

\subsection{Reproduction Command and Artifact Checks}

The complete formal design can be regenerated from the project root with
the following command:

\begingroup\small
\begin{verbatim}
python compare_decentralized_centralized.py `
  --T 1000 --nsim 20 `
  --seed-root 42 --seed-namespace FORMAL `
  --stop-alive-threshold 1 `
  --den-matcher-mode load --cen-matcher-mode off `
  --all-scenarios --force-rerun
\end{verbatim}
\endgroup

The command runs the four rollover/support scenarios and reconstructs the
CAR-threshold and weight measurement sweeps from their recorded state
paths. The separate loan-cap robustness run uses 20 paired replications
at every cap; the \(B=1200\) network table reuses the ON/ON artifacts.
For figure-only regeneration, \texttt{--skip-run} is used after artifact
validation. Each
artifact stores the schema version, code fingerprint, realized seeds,
matching mode, GNN checkpoint identity, feature switches, and observed
run lengths. The comparison program rejects incompatible or incomplete
artifacts instead of silently combining results from different code
versions. Calibration uses the separate \texttt{CALIBRATION} seed
namespace and is not included in the formal inference sample.

The public source archive for this study is available from the
\begin{center}\small
\href{https://github.com/zhanghaoyu1232-debug/central-bank-simulation}{\texttt{GitHub repository}}.
\end{center}
The executable source version used for the reported results is published on
the repository's \texttt{master} branch at Git commit
\begin{center}\small
\href{https://github.com/zhanghaoyu1232-debug/central-bank-simulation/commit/fe3445374adac3e90c2d29cbf90100d02b6ee135}{\nolinkurl{fe3445374adac3e90c2d29cbf90100d02b6ee135}}.
\end{center}
The repository and this commit were accessed and verified on September 8,
2026.
The versioned raw-output archive, \texttt{figures.zip}, is distributed with
the
\begin{center}\small
\href{https://github.com/zhanghaoyu1232-debug/central-bank-simulation/releases/tag/thesis-2026-09-08}{\texttt{Archived thesis release (thesis-2026-09-08)}}.
\end{center}
The annotated release tag resolves to the same executable-source commit stated
above.
Its core-source fingerprint is \texttt{11dda58061ded609}, matching the
recorded formal artifacts. The fingerprint identifies the nine core source
files used by the artifact compatibility check, whereas the Git commit
identifies the published repository snapshot. The supplementary reproduction manifest
records the full source-file checksums, trained-checkpoint checksum,
artifact schema, and realized seed sets. The accompanying README specifies
the execution commands, dependency requirements, raw-output paths, and
the scripts for regenerating the revised figures.

The numerical results in Chapters~6--7 are generated by this workflow.
The computational environment is Windows 11,
Python 3.10.11, NumPy 1.26.4, and Matplotlib 3.10.1. Formal paired seeds
are generated from \texttt{seed-namespace FORMAL} and
\texttt{seed-root 42}; the 20 realized ON/ON seeds are
\begin{quote}\small\raggedright
142333896, 189077467, 235919858, 283194794, 576266650, 614242082,
730482077, 742397554, 1076140681, 1193124925, 1558245143, 1943491583,
2598636728, 3288100970, 3356860510, 3363951558, 3381681453, 3647799003,
3716205366, and 4207354830.
\end{quote}
Raw replication files and figures are written under the project's versioned
output directory and archived as \texttt{figures.zip} in the release linked
above. The complete executable source, rather than the condensed Appendix~A
listings, is the electronic supplement for reproduction.

The loan-cap sweep uses a separate 20-seed sample, shared by both mechanisms
and all three cap values. In sorted order, these seeds are
\begin{quote}\small\raggedright
369133708, 383329927, 404488628, 550243861, 865305066,
1859786275, 1884968544, 1934392872, 2204291346, 2261209995,
2811363264, 2995172877, 3081541439, 3160032368, 3269070126,
3324115916, 3376120482, 3606691181, 3687649985, 4190266092.
\end{quote}
None of these seeds belongs to the formal ON/ON baseline sample. Thus the
\(B=1200\) sweep cell is an independent replication of the baseline
parameter setting, not a reuse of its numerical observations.

With 30 banks, a single run is computationally modest; the larger cost
comes from repeating stochastic paths and scanning parameters. This
scale was chosen so that contract histories and network snapshots remain
interpretable while still allowing heterogeneous bank types. It should
not be read as a claim that a 30-node network reproduces the size
distribution of an actual national banking system.

\subsection{Basic Simulation Parameters}

The baseline experiments use a total population of
\(N=30\) institutions. The institutional composition and type-specific
settings are reported in Section~6.2. This subsection summarizes the
parameters applied to both centralized matching and decentralized RFQ.

One simulation period represents one business day. The formal experiment
uses the pre-specified horizon
\begin{equation}
T=1000.
\end{equation}
A run stops earlier only when the number of surviving non-central banks is
one or fewer. The central bank is excluded from this count. Network
stability is retained as a diagnostic variable but does not terminate the
formal baseline or sensitivity runs. If collapse is not observed by
\(T\), the collapse time is right-censored at the horizon; risk-path and
bank-day statistics remain available over the full recorded interval.

The baseline experiments include bull and bear market states. Their
effects on interest rates, investment returns, liquidity, funding
outflows, and bank behavior follow the market-update rules defined in
Section~3.3.

The main parameter settings are summarized in
Table~\ref{tab:baseline-parameters}.

\begin{longtable}{@{}>{\raggedright\arraybackslash}p{3.25cm}>{\centering\arraybackslash}p{1.65cm}>{\centering\arraybackslash}p{3.55cm}>{\raggedright\arraybackslash}p{4.90cm}@{}}
\caption{Baseline simulation parameters}\label{tab:baseline-parameters}\\
\toprule
\textbf{Parameter} & \textbf{Symbol} & \textbf{Baseline setting} & \textbf{Description}\\
\midrule
\endfirsthead
\multicolumn{4}{c}{\tablename\ \thetable\ continued}\\
\toprule
\textbf{Parameter} & \textbf{Symbol} & \textbf{Baseline setting} & \textbf{Description}\\
\midrule
\endhead
\bottomrule
\endfoot
Number of institutions & $N$ & 30 & Simulated population\\
Formal horizon & $T$ & 1000 business days & Pre-specified evaluation horizon\\
Early-stop threshold & -- & 1 surviving non-central bank & Operational systemic-collapse threshold; central bank excluded\\
Market states & -- & Bull / Bear & Exogenous market regimes\\
CAR threshold & $\theta$ & 0.08 & Capital-stress threshold\\
Single-loan cap & $\bar x$ & 1200 & Maximum bilateral loan size\\
Monte Carlo repetitions & $M$ & 20 & Replications per baseline setting\\
Systemic-risk weights & $\mathbf w$ & $(0.5,\,0.3,\,0.2)$ & Weights for $\mathrm{FR}$, $\mathrm{CBS}$, and $\mathrm{CGR}$\\
Centralized matching cycle & -- & 1 business day & Market-wide rate-based allocation each day\\
RFQ rounds & -- & 4 & Maximum quotation rounds per business day\\
RFQ candidates per round & $K$ & 4 & Maximum lenders contacted from the local pool in each round\\
RFQ local discovery budget & -- & 8 & New local lenders per borrower and day; outstanding lenders excluded\\
RFQ random exploration & -- & At most 1 lender & Exploration within the discovery budget\\
RFQ degree cap & -- & 12 & Maximum local network degree for a new pair\\
Minimum trade size & -- & $10^{-9}$ & Baseline RFQ floor; effectively inactive\\
OFF contract maturity & -- & 1 business day & Next-day single payment of principal and interest\\
ON contract tenor & -- & 5--20 business days & Amortizing schedule selected at trade origination\\
Repayment structure & -- & OFF: single payment; ON: installment & Rollover switch selects the original contract schedule\\
Rollover mechanism & -- & Enabled & Active in the baseline configuration\\
Policy support & -- & Enabled & Liquidity and solvency/capital support are available\\
Liquidity-support total budget & -- & 12000 & Simulation-unit aggregate liquidity-support budget\\
Liquidity-support step budget & -- & 800 & Maximum liquidity-support budget reset each business day\\
Solvency-support total budget & -- & 6000 & Simulation-unit aggregate solvency-support budget\\
Solvency-support step budget & -- & 400 & Maximum solvency-support budget reset each business day\\
Common-shock probability & -- & 0.035 per day & Equity-neutral deposit run plus liquidity and credit stress\\
Common-shock outflow multiplier & -- & 1.25 & Bank-specific same-day outflow pressure, capped at 0.95\\
Common-shock project-PD multiplier & -- & 2.0 for 5 days & Temporary credit-risk persistence after a common shock\\
Initialization failure draw & -- & 0.03 at day 0 & Conditional on the draw, one non-central bank is selected uniformly\\
Random-seed base & $s_0$ & 42 & Base seed used to generate replication seeds\\
\end{longtable}

The baseline calibration enables both rollover and policy support. When
rollover is enabled, each new interbank loan is originated directly as a
5--20-period amortizing installment contract. When rollover is disabled,
it is originated as a one-business-day single-payment contract. The
current rule does not use a fixed rollover fraction and does not convert a
matured bullet contract ex post.

Policy support contains both liquidity-like and solvency/capital-like
instruments. The policy rule observes the systemic-risk value recorded in
the preceding period, adjusts the policy environment, and may allocate
support subject to per-bank and per-step budget constraints. Liquidity
support identified after the current Eisenberg--Noe settlement is queued
and disbursed at the start of the next business day; there is no
contemporaneous pre-clearing cash backstop inside the Eisenberg--Noe fixed
point. Solvency/capital support is applied after project returns have been
realized and before the final negative-equity cascade of the day. This
timing keeps liquidity support, rollover, and capital repair as distinct
mechanisms.

The capital adequacy ratio threshold \(\theta\) determines capital
stress. In the baseline setting,
\begin{equation}
\theta=0.08.
\end{equation}
Later sensitivity analysis examines how alternative threshold values
change the classification and measured systemic-risk path.

The single-loan cap \(\bar{x}\) restricts the size of each new bilateral
interbank transaction:
\begin{equation}
0\leq x_{ij,t}\leq\bar{x}.
\label{eq:loan-cap}
\end{equation}
The same cap is applied to both mechanisms.

The base random seed is \(s_0=42\). Named seed streams separate formal
experiments from calibration and training runs. Within the formal stream,
the same realized replication seed is supplied to centralized and RFQ
simulations. The realized seeds and experiment namespace are written to
the exported artifacts, and pairing is verified from these stored values
before mechanism-level statistics are calculated.

The default systemic-risk weights are
\begin{equation}
\left(w_1,w_2,w_3\right)=\left(0.5,0.3,0.2\right),
\end{equation}
\vspace{-0.4\baselineskip}
\noindent
where \(w_1\) weights the failure rate, \(w_2\) the CAR-breach share,
and \(w_3\) the capital-gap ratio. These are selected evaluation weights
rather than empirically estimated regulatory parameters.

\subsection{Experimental Design}

The experiment compares two complete interbank transaction designs under
common bank populations, balance-sheet initialization rules, regulatory
parameters, contract accounting, and due-payment clearing. Within each
paired Monte Carlo replication, the two mechanisms use the same realized
replication seed so that the initial conditions and exogenous stochastic
environment are comparable.

Both mechanisms operate every business day. Centralized matching observes
all eligible lenders and borrowers at once and requires complete borrower
funding before tentative trades are committed. Decentralized RFQ restricts
each borrower to a locally constructed information pool, executes bilateral
quotes sequentially, and retains acceptable partial trades. The resulting
Centralized--RFQ difference is therefore a comparison of market-wide
coordination with local search, screening, and bilateral execution; it is
not confounded by an imposed difference in trading frequency.

Within each simulation day, queued liquidity support and maturing
central-bank support contracts are processed first. Scheduled interbank
cash flows are then constructed from the contract book and settled through
one gross Eisenberg--Noe system. Defaults and creditor write-offs are
resolved to a fixed point before matching. After market conditions and bank
states are refreshed, the applicable daily matching mechanism forms new
contracts. Project returns are then realized, solvency support is applied
when enabled, and negative-equity defaults and their write-offs are again
resolved to a fixed point. Newly originated trades enter the outstanding-
exposure network immediately but enter clearing only on their scheduled
payment dates.

Formal simulations use the pre-specified horizon \(T=1000\). A run ends
earlier only if one or fewer non-central banks remain active. Network
stability is recorded as a diagnostic and is not a stopping rule. For the
rare early-collapse pair with realized lengths \(T_a\) and \(T_b\), the
common comparison horizon is
\begin{equation}
H_{a,b}=\min\{T_a,T_b\}.
\label{eq:common-horizon}
\end{equation}
Endpoint differences and integrated-risk measures are aligned within this
overlap. When neither run collapses early, \(H_{a,b}=1000\). Collapse time,
restricted mean survival time, and bank-days are reported separately so
that right-censored survival information is not inferred from padded risk
paths.

\section{Bank Types}

The definitions and behavioral characteristics of the three institution
types are presented in Section~3.2. This section reports the numerical
composition of the simulated banking system and the type-specific
settings used in the baseline experiments.

The simulated system contains 30 institutions: one central bank, 20
commercial banks, and nine shadow banks. Bank 0 represents the central
bank, banks 1--20 are commercial banks, and banks 21--29 are shadow
banks. Commercial and shadow banks constitute the 29 ordinary
institutions that may participate in centralized or RFQ matching.

The type-specific experimental settings are summarized in Table~\ref{tab:bank-types}.


\begin{table}[htbp]
\centering
\caption{Type-specific experimental settings}
\label{tab:bank-types}
\scriptsize
\setlength{\tabcolsep}{3.2pt}
\renewcommand{\arraystretch}{1.18}
\begin{tabularx}{\textwidth}{@{}>{\raggedright\arraybackslash}p{1.55cm}>{\centering\arraybackslash}p{0.85cm}>{\raggedright\arraybackslash}p{4.30cm}>{\raggedright\arraybackslash}p{2.65cm}X@{}}
\toprule
\textbf{Bank type} & \textbf{No.} & \textbf{Balance-sheet and funding characteristics} & \textbf{Risk-appetite treatment} & \textbf{Role in the model}\\
\midrule
Central bank & 1 & Larger capital and liquidity buffers; policy-related funding & Not applicable to ordinary matching & Sets the policy-rate environment and provides liquidity support when enabled\\
Commercial banks & 20 & Type-specific capital, liquidity, and liability ranges; deposits, interbank liabilities, and wholesale funding & Updated according to market conditions and bank states & Main ordinary participants in interbank lending\\
Shadow banks & 9 & More heterogeneous balance sheets; interbank and wholesale funding oriented; deposit ratio set to zero & Same baseline rule as commercial banks unless separately calibrated & Introduce heterogeneity in funding structure and refinancing dependence\\
Total & 30 & -- & -- & Complete simulated banking system\\
\bottomrule
\end{tabularx}
\end{table}


The type classifications determine the initial balance-sheet and
funding-structure settings but do not fix an institution's daily market
role or network position. Commercial and shadow banks may become
lenders, borrowers, or non-participants according to their updated
financial states, as described in Section~3.2.

\section{Experimental Scenarios}

Controlled comparisons organize the experimental design. Bank
population, initialization rules, regulatory parameters, contract
accounting, and clearing procedures remain unchanged across paired
simulations. The baseline enables both rollover and policy support, and
the main feature-ablation experiment evaluates all four combinations of
these two switches.

\subsection{Baseline Scenario}

The baseline financial system contains 30 institutions: one central
bank, 20 commercial banks, and nine shadow banks. One simulation period
represents one business day. The formal evaluation horizon is
\(T=1000\). A run terminates before this date only when the number of
active non-central banks falls to one or fewer; otherwise it is treated
as right-censored at day 1000 for collapse-time analysis.

Capital stress is evaluated using
\begin{equation}
\theta=0.08,
\end{equation}
and composite systemic risk uses
\begin{equation}
\left(w_1,w_2,w_3\right)=\left(0.5,0.3,0.2\right).
\end{equation}
Both rollover and policy support are enabled in the baseline. Daily
records include
\begin{equation}
\mathrm{FR}_t,\qquad\mathrm{CBS}_t,\qquad\mathrm{CGR}_t,\qquad\mathrm{SR}_t,
\end{equation}
as well as matching, clearing, network, and policy-support metrics.

\subsection{Centralized Matching Scenario}

Centralized matching provides a daily high-coordination benchmark. On
each business day, all eligible lenders and borrowers are considered
jointly subject to available supply, demand, reservation rates, the
single-loan cap \(\bar{x}\), and degree constraints.

An all-or-nothing rule applies at the borrower level. Funding is
committed only when compatible lenders can cover the borrower's complete
demand; otherwise, tentative allocations are cancelled. Consequently,
centralized coordination provides broad counterparty access but may leave
an otherwise feasible borrower entirely unfunded when complete coverage
cannot be assembled.

\subsection{Decentralized RFQ Scenario}

Decentralized RFQ represents daily local enquiry and bilateral quotation.
Each borrower first constructs one local information pool. Outstanding
lenders needed for rollover are considered first and do not consume the
new-discovery budget. The remaining pool is assembled from historical
counterparties, initial acquaintances and current one-hop exposure
neighbors, two-hop referrals, and at most one random exploratory lender.
At most eight newly discovered lenders enter the daily pool. The borrower
may then conduct up to four RFQ rounds and contact at most \(K=4\) lenders
per round. Lenders quote according to their remaining cash, bilateral
room, reservation rates, borrower risk, and post-trade concentration.

Unlike centralized allocation, RFQ retains acceptable partial trades even
when they do not cover the borrower's complete demand. Limited search,
rejected quotations, and insufficient candidate-lender supply can
therefore leave positive unmet demand. The baseline minimum trade size
\(m=10^{-9}\) is not a material screening constraint.

\subsection{Rollover and Policy-Support Scenarios}

The principal feature comparison uses the four combinations
\begin{equation}
(\mathrm{Rollover},\mathrm{Support})
\in
\{(\mathrm{ON},\mathrm{ON}),
(\mathrm{ON},\mathrm{OFF}),
(\mathrm{OFF},\mathrm{ON}),
(\mathrm{OFF},\mathrm{OFF})\}.
\end{equation}
Rollover changes contractual timing at origination: ON creates a
5--20-period installment schedule, whereas OFF creates a next-day
single-payment obligation.
Policy support changes the availability of next-period liquidity
assistance and same-day post-project solvency/capital assistance. These
switches are applied symmetrically to both matching mechanisms.

Scenario comparisons use the common pre-specified 1000-day horizon when
runs remain above the collapse threshold. If a run collapses earlier,
path comparisons use the paired observed overlap, while survival outcomes
are handled with right-censoring-aware statistics rather than treating
missing post-collapse observations as recoveries.

\subsection{Measurement and Structural Sensitivity Scenarios}

Measurement sensitivity examines the CAR threshold \(\theta\) and the
weights \((w_1,w_2,w_3)\) of the composite systemic-risk indicator.
Changing \(\theta\) reclassifies the recorded CAR paths and therefore
changes \(\mathrm{CBS}_t\), \(\mathrm{CGR}_t\), and the measured
\(\mathrm{SR}_t\). Changing the weights recombines the same recorded
component paths. In the reported sensitivity figures, neither operation
reruns the financial system or changes the underlying bank states.

Alternative weight combinations satisfy
\begin{equation}
w_1+w_2+w_3=1,\qquad w_1,w_2,w_3\geq0.
\end{equation}
The component paths \(\mathrm{FR}_t\), \(\mathrm{CBS}_t\), and
\(\mathrm{CGR}_t\) are therefore retained alongside \(\mathrm{SR}_t\).

As a structural robustness check, the single-loan cap is varied over
\begin{equation}
\bar{x}\in\{600,1200,1800\},
\end{equation}
\vspace{-0.4\baselineskip}
\noindent
with \(\bar{x}=1200\) as the baseline. This changes the feasible size of
a bilateral loan and can therefore alter funding allocation, exposure
concentration, and network structure rather than merely rescaling the
risk indicator.

\subsection{Summary of Experimental Scenarios}

\begin{table}[htbp]
\centering
\caption{Experimental scenarios and principal outcomes}
\label{tab:experimental-scenarios}
\small
\renewcommand{\arraystretch}{1.16}
\begin{tabularx}{\textwidth}{@{}>{\raggedright\arraybackslash}p{0.22\textwidth}>{\raggedright\arraybackslash}p{0.34\textwidth}X@{}}
\toprule
\textbf{Scenario} & \textbf{Changed element} & \textbf{Main outcomes}\\
\midrule
Baseline & Rollover ON, support ON & Reference paths, timing, and network evolution\\
Centralized matching & Daily market-wide allocation with a borrower-level full-funding rule & Exposure distribution, funding allocation, and clearing-related risk\\
Decentralized RFQ & Daily local enquiries, four rounds, $K=4$, discovery budget 8, partial funding & Bilateral network formation, screening, funding outcomes, and clearing-related risk\\
Feature ablation & Rollover/support ON--OFF combinations & Timing of systemic risk and mechanism--policy interaction\\
CAR sensitivity & $\theta$ & Capital-stress classification and measured risk path\\
Weight sensitivity & $(w_1,w_2,w_3)$ & Dependence of $\mathrm{SR}_t$ on measurement choices\\
Loan-cap robustness & $\bar{x}\in\{600,1200,1800\}$ & Composite-risk ranking; network and funding statistics at the baseline cap\\
\bottomrule
\end{tabularx}
\end{table}

Synthetic balance sheets, stylized market regimes, daily centralized
coordination, and a 30-institution population support
controlled computational comparisons but limit external validity. The
reported numerical levels should therefore be interpreted as model
outcomes rather than forecasts for a particular banking system.

\section{Simulation Procedure}

The conceptual timing structure is presented in Section~3.5. This
section specifies the executable simulation procedure used for
initialization, daily iteration, output recording, Monte Carlo
aggregation, parameter sweeps, and simulation-horizon control.

\subsection{Initialization}

For each Monte Carlo replication, the simulation performs the following
initialization procedure:

\begin{enumerate}
\def\labelenumi{\arabic{enumi}.}
\item
  Create the bank population using the institutional composition and
  type-specific settings reported in Section~6.2.
\item
  Initialize core capital, liquid assets, current liabilities, project
  assets, funding composition, risk appetite, and funding outflow rates
  according to the applicable bank-type rules.
\item
  Calculate the initial capital adequacy ratio, liquidity coverage
  ratio, solvency ratio, and model-specific leverage measure from the
  initialized balance sheets.
\item
  Initialize the market state and the corresponding policy rates,
  lending rates, investment returns, volatility, and funding conditions
  according to the mechanism defined in Section~3.3.
\item
  Create an empty contract book and initialize the due-payment liability
  matrix and outstanding-exposure matrix.
\item
  Set the matching mechanism, formal horizon, early-collapse threshold,
  CAR threshold, single-loan cap, named seed stream, and other
  parameters reported in Section~6.1.2.
\item
  Initialize the time-series containers used to record bank states,
  clearing results, matching outcomes, network measures, and
  systemic-risk indicators.
\end{enumerate}

\subsection{Multi-Period Simulation Process}

For each business day, the simulation executes the following sequence:

\begin{enumerate}
\def\labelenumi{\arabic{enumi}.}
\item
  Disburse any liquidity support queued after the preceding day's
  settlement. Reset the daily policy budgets, observe the previously
  recorded systemic-risk state, and update the central-bank policy rule.
\item
  Process repayments of previously issued central-bank support loans
  that mature in the current period.
\item
  Extract all single-payment and installment amounts currently due from
  the interbank contract book, including previously accumulated arrears.
\item
  Aggregate all due contracts into one gross debtor-to-creditor
  due-payment matrix and solve one Eisenberg--Noe fixed point. No
  contemporaneous liquidity facility is inserted into this clearing
  matrix.
\item
  Apply Eisenberg--Noe clearing, allocate each debtor's payment pro rata
  across its due contracts, and update or remove contracts. Under
  rollover OFF, a residual single-payment balance causes hard default;
  under rollover ON, an unpaid installment becomes arrears and does not
  cause default solely because of the missed payment.
\item
  Rebuild interbank asset and liability stocks, refresh regulatory
  measures, resolve creditor write-offs and any induced defaults to a
  fixed point, and queue eligible liquidity support for the next
  business day.
\item
  Update the market state, policy-related rates, investment returns,
  volatility, funding outflows, liquid assets, current liabilities, and
  exogenous shocks.
\item
  Assign each active non-central bank as a lender, borrower, or
  non-participant using its updated state.
\item
  Apply the selected transaction design. Under centralized matching,
  all current intentions enter the daily market-wide allocation. Under
  decentralized RFQ, quotations are requested sequentially from the
  borrower's current local information pool.
\item
  Register completed trades as new one-period single-payment contracts
  when rollover is OFF or 5--20-period installment contracts when it is
  ON. Allocate borrowed funds according to the project/reserve rules and
  update the project book.
\item
  Realize project returns, refresh accounting and regulatory measures,
  and apply eligible solvency/capital support. Classify negative-equity
  failures and resolve their creditor write-offs recursively to a fixed
  point.
\item
  Reconstruct outstanding interbank exposures and calculate
  \(\mathrm{FR}_t\), \(\mathrm{CBS}_t\), \(\mathrm{CGR}_t\), and
  \(\mathrm{SR}_t\).
\item
  Update the network-stability diagnostic. Terminate only if one or
  fewer non-central banks remain active; otherwise continue through the
  pre-specified day \(T=1000\).
\end{enumerate}

The end-of-period bank states, active contracts, and outstanding
exposures are carried forward to the next business day.

\subsection{Recording of Risk Indicators}

End-of-period records contain bank states, active contracts, outstanding
exposures, clearing outcomes, matching outcomes, policy-support amounts,
and systemic-risk indicators. Failure status follows the absorbing
inactivity rule specified in Section~5.2: once a non-central bank fails
because of an OFF-regime unpaid single-payment balance or negative
accounting equity, it
remains inactive and continues to be counted in \(\mathrm{FR}_t\).

Funding records include
\begin{equation}
\mathrm{TotalSupply}_t,\quad
\mathrm{TotalDemand}_t,\quad
\mathrm{ActualLent}_t,\quad
\mathrm{FundingRatio}_t,\quad
\mathrm{UnmetDemand}_t.
\end{equation}
Funding satisfaction is
\begin{equation}
\mathrm{FundingRatio}_t
=
\frac{\mathrm{ActualLent}_t}
{\max\{\varepsilon,\mathrm{TotalDemand}_t\}},
\end{equation}
and unmet demand is
\begin{equation}
\mathrm{UnmetDemand}_t
=
\max\{0,\mathrm{TotalDemand}_t-\mathrm{ActualLent}_t\}.
\end{equation}
Because both mechanisms match daily, two common funding measures are
reported. The first is the event-conditional ratio
\(\mathrm{FundingRatio}_t\), which uses only business days with positive
borrower demand. The second is cumulative transaction volume, which sums
\(\mathrm{ActualLent}_t\) over the common calendar horizon, including days
with little or no matching activity. Unmet demand is the corresponding
calendar-day residual \(\mathrm{UnmetDemand}_t\) defined above.
Days without positive borrower demand are excluded from the ratio rather
than being coded as either perfect or zero funding.

Network statistics are calculated from the gross outstanding
creditor--debtor exposure network, excluding the central bank from
ordinary interbank-network measures. Recorded variables include active
directed links, density, degree and weighted-degree moments, the largest
bilateral exposure, exposure HHI, the largest weakly connected component,
repeated counterparties, transaction volume, funding satisfaction,
unmet demand, and clearing shortfall.

The artifact pipeline stores time-averaged structural measures within
each run and cumulative transaction/shortfall measures for that run.
Under the formal 1000-day design, these measures share the same horizon
unless early systemic collapse occurs. If a pair ends at different dates,
cumulative quantities are compared on the pairwise overlap or normalized
by observed business days rather than ranked by unmatched terminal totals.

Contract-book records preserve lender--borrower identities, remaining
principal, interest rates, and future payment dates. Single-path diagnostic
logs can record dated policy issuance events, but the exported Monte Carlo
artifacts store only cross-seed mean daily liquidity and solvency support
paths. Chapter~7 therefore compares intervention intensity from those
series over the first 100 business days and over the full 1000-day horizon:
mean support per day, the number of days with positive support in the short
window, and the first positive-support date. Unique-recipient counts,
repayment rates, and remaining budgets are not estimated, because those
fields are not stored in the formal artifacts.

\subsection{Monte Carlo Repetitions}

The baseline trajectories use \(M=20\) Monte Carlo replications. Each
mechanism records the realized seed of every run, and Centralized--RFQ
statistics are calculated only for seeds present in both artifacts. This
provides an explicit paired-replication check rather than assuming that
two runs are paired solely because they occupy the same list position.

Let \(T_m^{(k,s)}\) denote the realized run length of replication
\(m\) under mechanism \(k\) and feature scenario \(s\). The exported
aggregate trajectory is reported through the pre-specified 1000-day
horizon when the run does not collapse early. For plotting only, an
early-collapsed path is extended by its last observed state. This
absorbing fill prevents \(\mathrm{FR}_t\), \(\mathrm{CBS}_t\), and
\(\mathrm{SR}_t\) from falling mechanically when a failed replication
leaves the contributor set. Filled terminal values are not used to infer
survival duration: formal survival outcomes use the recorded collapse
step, restricted mean survival time, and non-central-bank days alive.

For descriptive mechanism-specific trajectories, pointwise 95\%
confidence intervals use the normal critical value implemented by the program,
\begin{equation}
\overline{\mathrm{SR}}_t
\pm
1.96
\frac{s_{\mathrm{SR},t}}{\sqrt{M_t}},
\label{eq:mc-ci}
\end{equation}
\vspace{-0.4\baselineskip}
\noindent
where \(M_t\) is the number of replications represented at business day
\(t\); after absorbing display fill it remains 20 for the plotted mean.
These intervals are pointwise rather than simultaneous confidence bands.
Component intervals are constructed in the same way.

For formal Centralized--RFQ path, endpoint, network, and funding
comparisons, the program first calculates one difference per matched seed
and reports the paired mean plus or minus (1.96) paired standard errors.
This keeps the unit of inference at the replication rather than treating
business days as independent observations. Bank-day and conditional
first-passage contrasts use a paired nonparametric bootstrap with 2,000
replicates and fixed seed 0, matching the executable implementation.

For a paired Centralized--RFQ replication with common seed \(m\), define
\begin{equation}
H_m=\min\left\{T_m^{\mathrm{RFQ}},T_m^{\mathrm{Centralized}}\right\}.
\label{eq:pair-native-horizon}
\end{equation}
The paired path difference is defined on that pair's common native window:
\begin{equation}
\Delta\mathrm{SR}_t^{(m)}
=\mathrm{SR}_t^{(m),\mathrm{RFQ}}
-\mathrm{SR}_t^{(m),\mathrm{Centralized}},\qquad 1\leq t\leq H_m.
\label{eq:paired-sr-difference}
\end{equation}
For scenario \(s\), the reported mean path contrast is
\begin{equation}
\overline{\Delta\mathrm{SR}}_s
=\frac{1}{\lvert\mathcal P_s\rvert}
\sum_{m\in\mathcal P_s}
\left(\frac{1}{H_m}\sum_{t=1}^{H_m}\Delta\mathrm{SR}_t^{(m)}\right),
\label{eq:paired-mean-path}
\end{equation}
where \(\mathcal P_s\) is the set of matched seeds. Each replication pair
receives equal weight after its within-pair time average is calculated.
The mean and uncertainty are computed across these paired summaries.
Endpoint contrasts use day \(H_m\); AUC contrasts integrate the two native
paths over the same pair-specific interval. They are not calculated from
last-state padding or from a changing set of available replications.

For the principal ON/ON sample, every \(H_m=1000\), so all path, endpoint,
and integrated-risk comparisons use the same fixed 1000-day interval.
All ON/OFF pairs also reach day 1000. In the OFF/ON and OFF/OFF samples,
the pair-specific horizons have means 994.7 and 732.45 days and minima
894 and 18 days, respectively. Their path contrasts therefore compare
matched observed lifetimes, rather than one fixed period common to every
pair. Recorded survival outcomes are reported separately to describe the
differences in observation length.

Timing is also treated as an outcome. The first-passage results reported
for ON/ON use the common fixed interval \([1,1000]\). For a selected
systemic-risk threshold \(q\), the hit rate is the share of replications
that reach \(q\) within that interval. The paired timing difference
\(\Delta\tau_q=\tau_q^{\mathrm{RFQ}}-\tau_q^{\mathrm{Centralized}}\)
is summarized only for pairs in which both mechanisms reach the threshold.
Its uncertainty uses the paired-bootstrap procedure described above.
Mechanism-specific hit rates are reported alongside the conditional timing
difference so that non-hits remain explicit.

\subsection{Sweep Experiments and Simulation Horizon}

The reported measurement sweeps use the same 20 paired 1000-day baseline
paths. They change only the mapping from recorded bank states to the
composite indicator; they do not create additional economic simulations.

For the CAR-threshold sweep, \(N_\theta=8\) values are examined from
0.08 to 0.15 around the baseline \(\theta=0.08\). For each value, the
stored bank-level CAR paths are reclassified, \(\mathrm{CBS}_t\) and
\(\mathrm{CGR}_t\) are recalculated, and the systemic-risk series is
reconstructed. This is a measurement-threshold sensitivity analysis,
not a policy-follow counterfactual.

The weight-sensitivity experiment uses \(w_1\in\{0.1,\ldots,0.9\}\) while preserving
\(w_2:w_3=3:2\) and \(w_1+w_2+w_3=1\). The weight-sensitivity experiment
rescales the recorded \(\mathrm{FR}\),
\(\mathrm{CBS}\), and \(\mathrm{CGR}\) paths rather than rerunning the
financial system for every weight combination.

A separate structural robustness experiment varies
\begin{equation}
\bar{x}\in\{600,1200,1800\}.
\end{equation}
The robustness experiment uses 20 paired replications at each cap. The
baseline network and funding table continues to use the formal ON/ON
sample at \(\bar{x}=1200\); the cap sweep is used to assess the direction
of the composite-risk ranking and is not interpreted as uniform
statistical separation.

\subsection{Figure and Table Outputs}

Recorded simulation histories produce time-series figures, network
snapshots, sensitivity plots, and summary tables. Mean systemic-risk
paths compare the two matching mechanisms through

\begin{equation}
\overline{\mathrm{SR}}_{t}.
\end{equation}

Risk-decomposition figures report

\begin{equation}
\overline{\mathrm{FR}}_{t},
\qquad
\overline{\mathrm{CBS}}_{t},
\qquad
\overline{\mathrm{CGR}}_{t}.
\end{equation}

Sensitivity plots show the complete family of curves using a continuous
colour scale, with the baseline \(\theta=0.08\) or \(w_1=0.5\) highlighted.
The loan-cap figure is reported separately because it represents a
structural rerun rather than a remeasurement of the same paths.

\begin{table}[htbp]
\centering
\caption{Simulation outputs}
\label{tab:simulation-outputs}
\small
\renewcommand{\arraystretch}{1.16}
\begin{tabularx}{\textwidth}{
@{}
>{\raggedright\arraybackslash}p{0.21\textwidth}
>{\raggedright\arraybackslash}p{0.39\textwidth}
X
@{}
}
\toprule
\textbf{Output type} & \textbf{Content} & \textbf{Purpose}\\
\midrule
Risk path figure
& $\overline{\mathrm{SR}}_{t}$ on common observed time
& Compare systemic-risk paths without terminal-horizon bias\\

Component figure
& $\overline{\mathrm{FR}}_{t}$,
  $\overline{\mathrm{CBS}}_{t}$,
  $\overline{\mathrm{CGR}}_{t}$
& Identify the source of changes in composite risk\\

Network snapshot
& Outstanding exposures and bilateral links
& Examine network structure and concentration\\

Sensitivity / robustness figure
& CAR-threshold, weight, or loan-cap sweep
& Assess measurement sensitivity and structural robustness\\

Policy-support path summary
& Cross-seed mean daily liquidity and solvency support
& Compare early-window and full-horizon intervention intensity\\

Run-length / first-passage table
& Observed $T_m$, collapse hits, and $\mathrm{SR}_t$ hitting times
& Report right-censoring and whether timing contrasts are identified\\
\bottomrule
\end{tabularx}
\end{table}
